\section{A common natural family for selected coefficients}

The following results provide supporting algebra and coarse counts for
\cite[Proposition~5.7]{Arad2016RigorousRG}.
The number of family members and their open-boundary bond dimensions are
separate quantities. No sharp source bond or running-time conclusion is
included.

\begin{theorem}[Natural middle operators chosen before the exterior pair]
    \label{thm:selected_common_natural_mpo_family}
    \lean{TNLean.MPS.AreaLaw.exists_common_physical_twoMarkerCoefficient_mpo_family}
    \uses{thm:selected_coefficient_prescribed_middle,
        thm:selected_middle_fixed_physical_data,thm:finite_natural_retained_core_family,
        thm:actual_natural_affine_product,thm:physical_interval_nested_trace,
        thm:physical_interval_retained_matrix,def:two_marker_middle_index}
    Let the chain have $L$ sites with local dimension $q\ge0$.
    Let $a,b,\ell$ be nonnegative integers satisfying
    $a+2\ell\le b$ and $b+2\ell\le L-1$, and choose $0\le p,s<2\ell$.
    Use the regions
    \begin{align}
        T & ={[0,a+p+1)}, & E &
        ={[a+p+1,b+s+1)},\notag                 \\
        R & =(T\cup E)^c, & Z & =[a+2\ell+1,b),
        \label{eq:selected_common_natural_regions}
    \end{align}
    with all intervals clipped to the chain. Set $N=|E|$, and identify $E$
    with its increasing list of sites. Let $\rho$ be the number of
    configurations on $E\setminus Z$.

    Suppose the supplied window operators $W^L_r,W^R_r$ are supported on
    $[a+r,a+r+2)$ and
    $[b+r,b+r+2)$, respectively.
    Supply site-dependent core tensors of virtual dimension $B$ on
    $[a+2\ell,b+1)$ and a closing matrix. Let $P_M$ be their trace contraction,
    canonically placed on the full chain.
    Use labels $L_r,R_r,B_0,B_1,B_2$, and put
    $O=\{L_r,R_r,B_0,B_1,B_2\}\setminus\{L_p,R_s\}$.
    Fix complex ordinary weights $z_t$ and nonnegative integers $k,i,j$.
    Write $\mathcal W_{k,i,j}$ for the two-marker middle index set with
    $q^2$ labels for each marker family, as in
    Definition~\ref{def:two_marker_middle_index}.

    There exist a finite family of natural-interval operators $V_r$,
    $0\le r<n$, and degrees $0\le d_r\le k$, such that
    \begin{align}
        n & \le |\mathcal W_{k,i,j}|(k+1)\rho^2,
        \label{eq:selected_common_natural_count}
    \end{align}
    and every $V_r$, on the paired physical alphabet, has an open-boundary
    MPS representation of bond at most $(B^{d_r})^2$ in increasing site order.
    These choices precede the following exterior data.

    For every supported exterior pair $P_L,P_R$ on $[0,a+1)$ and
    $[b+2\ell,L)$, respectively, form the full family
    $F_{L_r}=W^L_r$, $F_{R_r}=W^R_r$,
    $F_{B_0}=P_L$, $F_{B_1}=P_M$, $F_{B_2}=P_R$, and set
    \begin{align}
        D_z & =\sum_{t\in O}z_tF_t.
        \label{eq:selected_common_natural_ordinary}
    \end{align}
    For every complex scalar family $c_0,\ldots,c_k$, there are joint exterior
    operators $Q_r$ on $\mathcal H_T\otimes\mathcal H_R$ such that
    \begin{align}
        \left(\sum_{h=0}^k c_h\,[u^iv^j]
        (D_z+uW^L_p+vW^R_s)^h\right)^{\mathrm{nat}}
         & =\sum_{r=0}^{n-1}Q_r\otimes V_r.
        \label{eq:selected_common_natural_decomposition}
    \end{align}
    The superscript first groups the two exterior regions and then reads the
    middle in its increasing coordinates. The formal variables $u,v$ commute
    with the matrix coefficients. The two selected crossings are unweighted;
    only the ordinary sum carries $z$.
    Zero physical or virtual dimensions, degree zero, and impossible crossing
    frequencies are included. No factorization of each joint exterior $Q_r$
    into separate left and right operators is asserted.
\end{theorem}
\begin{proof}
    \uses{thm:selected_coefficient_prescribed_middle,
        thm:selected_middle_fixed_physical_data,thm:finite_natural_retained_core_family,
        thm:actual_natural_affine_product,thm:physical_interval_nested_trace,
        thm:physical_interval_retained_matrix,def:two_marker_middle_index,
        thm:two_marker_prescribed_span}
    The prescribed-middle coefficient decomposition first gives one ordered
    middle factor for every index $w\in\mathcal W_{k,i,j}$, before the exterior
    pair or the scalars $c_h$ are chosen. Its ordinary factor has the form
    $z_{B_1}M_E+K_z$, with $K_z$ supported on $E\setminus Z$; its two marker
    families are also supported there.

    Represent those supported factors on the single retained complement.
    Replace each ordinary run by its repeated affine steps, with scalar
    $z_{B_1}$ on each ordinary step and scalar zero on each marker step.
    The resulting ordered list has length $h_w\le k$ and evaluates to the
    same middle word. The canonical interval and retained-coordinate
    identities identify its central trace operator with the supplied core.
    The finite natural family theorem expresses that word as a sum of exactly
    $(h_w+1)\rho^2$ natural operators, each with its corresponding square bond.

    Take the disjoint union over $w$ and enumerate it. Its size is at most
    $|\mathcal W_{k,i,j}|(k+1)\rho^2$, proving
    (\ref{eq:selected_common_natural_count}). Every prescribed middle word lies
    in the span of this one family. The prescribed-spanning-family
    factorization then supplies the exterior coefficients in
    (\ref{eq:selected_common_natural_decomposition}), without changing the
    chosen middle operators.
\end{proof}

\begin{theorem}[A polynomial-base bound for the selected family count]
    \label{thm:selected_common_natural_cardinality}
    \lean{TNLean.MPS.AreaLaw.actual_common_physical_mpo_card_le}
    \uses{thm:selected_middle_collar_cardinality,thm:two_marker_middle_cardinality,
        def:physical_interval_complement_coordinates}
    Retain the geometric data and positions above, and suppose $q\ge1$ and
    $k\ge1$. For arbitrary nonnegative $i,j,n$, the numerical implication
    \begin{align}
        n\le |\mathcal W_{k,i,j}|(k+1)\rho^2
        \quad\Longrightarrow\quad
        n\le\bigl(q(k+1)\bigr)^{8\ell+2(i+j)+2}
        \label{eq:selected_common_natural_cardinality}
    \end{align}
    holds. No low-frequency restriction is required.
\end{theorem}
\begin{proof}
    \uses{thm:selected_middle_collar_cardinality,thm:two_marker_middle_cardinality,
        def:physical_interval_complement_coordinates}
    The collar bound gives $\rho\le q^{4\ell}$, and hence
    $\rho^2\le q^{8\ell}$. With $m=i+j$ and $v=k+1$, the middle index count is
    at most $2^m q^{2m}v^{m+1}$. Since $v\ge2$, it is at most
    $q^{2m}v^{2m+1}$. Multiplication by $v\rho^2$ gives
    $q^{8\ell+2m}v^{2m+2}$.
    Increasing both exponents to $8\ell+2m+2$ proves
    (\ref{eq:selected_common_natural_cardinality}).
\end{proof}

\begin{theorem}[The low-frequency numerical consequence]
    \label{thm:selected_common_natural_low_frequency_cardinality}
    \lean{TNLean.MPS.AreaLaw.actual_common_physical_mpo_card_low_frequency_le}
    Let $q\ge2$, $k\ge1$, $\ell\ge1$, and let $i,j,n$ be nonnegative integers
    with $i,j\le\lfloor k/(2\ell)\rfloor$. Then
    \begin{align}
        n\le\bigl(q(k+1)\bigr)^{8\ell+2(i+j)+2}
        \quad\Longrightarrow\quad
        n\le(qk)^{20\ell+4\lfloor k/\ell\rfloor}.
        \label{eq:selected_common_natural_low_frequency_cardinality}
    \end{align}
    This is an arithmetic implication and supplies no further operator or
    bond assertion.
\end{theorem}
\begin{proof}
    The frequency bounds give $(i+j)\ell\le k$, so
    $i+j\le\lfloor k/\ell\rfloor$. Since $\ell\ge1$,
    $8\ell+2(i+j)+2\le10\ell+2\lfloor k/\ell\rfloor$.
    Moreover $q(k+1)\le(qk)^2$. Substituting this base comparison doubles the
    last exponent and gives the stated bound.
\end{proof}
